% Capítulo 3: Metodologia
\chapter{Metodologia}\label{chp:metodologia}

Este capítulo apresenta a metodologia para investigação comparativa de 
técnicas de processamento de imagens hiperespectrais em arquiteturas 
heterogêneas. A abordagem estrutura-se em oito seções: considerações iniciais, 
estratégia de investigação, procedimentos de investigação, estudo comparativo 
de arquiteturas FPGA e GPU, protocolo de simulação e benchmarking, metodologia 
de síntese arquitetural, planejamento de execução, e considerações finais. A 
metodologia privilegia análise teórica e simulada sobre implementação física, 
justificada pela necessidade de comparação entre abordagens tecnológicas e 
pela complexidade de validação experimental em múltiplas plataformas 
heterogêneas.

\section{Considerações Iniciais}\label{sec:consideracoes_iniciais}

A metodologia baseou-se na análise do levantamento bibliográfico apresentado 
no \autoref{chp:levantamento}, que revelou lacunas na literatura sobre sistemas 
integrados de processamento hiperespectral. A revisão identificou que apenas 
parte dos trabalhos abordam sistemas completos integrados. Esta constatação 
direcionou a pesquisa para abordagem comparativa, priorizando investigação 
teórica e simulada sobre implementação física de protótipos.

A pesquisa comparativa fundamenta-se em três aspectos. Primeiro, a análise 
das abordagens tecnológicas disponíveis, considerando que trabalhos como 
\textcite{Diaz2019} demonstram processamento de mais de 330 frames por segundo 
em GPU embarcada NVIDIA Jetson TX2 com algoritmo HyperLCA para compressão 
hiperespectral, enquanto \textcite{Hwang2011} propõe metodologia de design 
consciente de energia para compressão de imagens hiperespectrais, evidenciando 
diversidade de soluções e ausência de critérios unificados de comparação. 
Segundo, a complexidade de implementar e validar múltiplas arquiteturas 
heterogêneas em ambiente experimental, considerando que cada plataforma (FPGA 
Xilinx, GPU NVIDIA, processadores ARM) requer ferramentas específicas, licenças 
proprietárias e expertise técnica. Terceiro, a contribuição reside na 
proposição de framework teórico unificado para comparação e seleção de 
técnicas, preenchendo lacuna metodológica na literatura.

A priorização da análise teórica e simulada baseia-se na natureza exploratória 
da pesquisa. A investigação visa estabelecer critérios comparativos entre 
abordagens arquiteturais, identificar trade-offs entre performance, consumo 
energético e precisão, e propor metodologia para seleção de técnicas para 
diferentes cenários de aplicação. Estes objetivos são alcançados através de 
modelagem matemática, simulação e análise comparativa ao invés de implementação 
física limitada a configurações específicas de hardware.

\section{Abordagem Metodológica}\label{sec:abordagem_metodologica}

A estratégia de pesquisa baseia-se em abordagem exploratória-comparativa com 
ênfase quantitativa para investigar técnicas e arquiteturas disponíveis para 
processamento de imagens hiperespectrais. Esta seção apresenta a caracterização 
da pesquisa, a estratégia de investigação comparativa e a delimitação do escopo 
tecnológico.

\subsection{Caracterização da Pesquisa}\label{subsec:caracterizacao_pesquisa}

A pesquisa caracteriza-se como exploratória-comparativa, baseada em simulação 
e modelagem teórica. O caráter exploratório justifica-se pela ausência de 
estudos que comparem arquiteturas heterogêneas para processamento 
hiperespectral, conforme evidenciado no levantamento bibliográfico. A dimensão 
comparativa emerge da necessidade de estabelecer critérios para avaliação de 
técnicas e arquiteturas, permitindo identificação de soluções para diferentes 
cenários de aplicação.

A abordagem quantitativa baseia-se em análise de desempenho teórico e 
modelado, utilizando métricas padronizadas para comparação entre técnicas. Os 
parâmetros incluem throughput (frames por segundo), latência (milissegundos 
por frame), consumo energético (watts), utilização de recursos (percentual de 
DSPs, BRAMs, LUTs), precisão de classificação (accuracy, F1-score) e 
eficiência computacional (operações por segundo por watt). A abordagem 
quantitativa permite comparação entre trabalhos que utilizam métricas diferentes 
ou cenários experimentais distintos.

O benchmarking de técnicas e arquiteturas constitui elemento central da 
metodologia, estabelecendo conjunto de testes e métricas para avaliação 
comparativa. Os benchmarks incluem datasets de referência (AVIRIS Indian 
Pines, Pavia University, Salinas Valley), cargas de trabalho de aplicações 
reais, e protocolos de medição que consideram variabilidade temporal e 
condições operacionais. Esta padronização permite comparação entre trabalhos da 
literatura e validação dos resultados teóricos propostos.

\subsection{Estratégia de Investigação Comparativa}
\label{subsec:estrategia_investigacao}

A metodologia de comparação estrutura-se em quatro etapas: mapeamento das 
técnicas disponíveis, catalogação de características e performance, análise 
comparativa baseada em critérios objetivos, e síntese de recomendações para 
diferentes cenários de aplicação. A metodologia de comparação garante cobertura 
do espaço de soluções e identificação de lacunas ou oportunidades de otimização.

Os critérios de seleção de técnicas e modelos baseiam-se em relevância para 
aplicações embarcadas, disponibilidade de dados de performance, 
representatividade no estado da arte, e potencial de integração em arquiteturas 
heterogêneas. Técnicas que atendem estes critérios são priorizadas para análise, 
enquanto abordagens com limitações ou dados insuficientes são documentadas mas 
não incluídas na comparação quantitativa.

O framework de avaliação de desempenho estabelece métricas e protocolos de 
medição que permitem comparação entre abordagens. As métricas incluem 
indicadores de performance (throughput, latência), eficiência (consumo 
energético, utilização de recursos), qualidade (precisão, robustez) e 
escalabilidade (suporte a diferentes resoluções e bandas espectrais). Os 
protocolos especificam condições experimentais, datasets de teste, e 
procedimentos de medição que garantem reprodutibilidade e comparabilidade dos 
resultados.

A análise multi-critério para otimização utiliza técnicas de tomada de decisão 
que consideram múltiplos objetivos conflitantes. A metodologia AHP (Analytic 
Hierarchy Process) é empregada para ponderação de critérios segundo cenários 
de aplicação, enquanto análise de Pareto identifica soluções não-dominadas no 
espaço de trade-offs. A análise multi-critério permite identificação de 
configurações para diferentes requisitos operacionais.

\subsection{Delimitação do Escopo Tecnológico}
\label{subsec:delimitacao_escopo}

O escopo tecnológico concentra-se em FPGA Xilinx das famílias Zynq e 
UltraScale+, justificado pela adoção em aplicações de processamento de sinais 
e disponibilidade de ferramentas de desenvolvimento. A escolha por FPGA Xilinx 
baseia-se na análise da literatura que demonstra uso desta plataforma em 
trabalhos de processamento hiperespectral embarcado, com destaque para 
\textcite{Hwang2011} que propõe metodologia de design consciente de energia 
para compressão de imagens hiperespectrais, e \textcite{ArucuIliev2025} que 
apresenta estudo comparativo sobre eficiência energética em processamento 
hiperespectral acelerado por FPGA.

As GPU embarcadas incluem plataformas NVIDIA Jetson (TX2, Xavier, Orin) e AMD 
equivalentes, selecionadas por representarem processamento paralelo embarcado 
com restrições energéticas. A análise de \textcite{Diaz2019} demonstra 
viabilidade de processamento de mais de 330 frames por segundo em Jetson TX2 
com algoritmo HyperLCA para compressão hiperespectral em tempo real, 
evidenciando o potencial destas plataformas para processamento hiperespectral 
embarcado. O trabalho de \textcite{Zhang2023} apresenta revisão sobre 
classificação de imagens hiperespectrais de sensoriamento remoto por UAV, 
destacando avanços em técnicas de deep learning para esta aplicação.

As técnicas de otimização abrangem paralelização em pipeline, otimizações de 
memória, técnicas de quantização, compressive sensing, seleção de bandas 
espectrais, e fusão de dados multi-sensor. Esta seleção baseia-se na análise 
da literatura que identifica as técnicas de otimização como promissoras para aplicações 
embarcadas, considerando trade-offs entre performance e consumo energético.

As limitações do estudo incluem foco em aplicações de sensoriamento remoto 
com UAV, consideração de datasets públicos para validação, análise limitada a 
faixa espectral de 400-2500 nm (visível e infravermelho próximo), e premissa 
de operação em tempo real com restrições energéticas inferiores a 20W. Estas 
limitações garantem foco da análise, evitando dispersão em cenários com 
requisitos conflitantes.

\section{Procedimentos de Investigação}\label{sec:procedimentos_investigacao}

Os métodos de pesquisa e análise comparativa estruturam-se em três etapas: 
levantamento e catalogação de técnicas disponíveis na literatura, modelagem 
teórica comparativa baseada em análise matemática de performance, e 
desenvolvimento de framework de simulação para validação dos modelos teóricos. Esta seção detalha os procedimentos de cada etapa, estabelecendo protocolos 
para garantir reprodutibilidade e validade dos resultados.

\subsection{Levantamento e Catalogação de Técnicas}
\label{subsec:levantamento_catalogacao}

O mapeamento de técnicas FPGA baseia-se em revisão da literatura com foco em 
trabalhos publicados entre 2018-2024 em conferências e periódicos (IEEE TGRS, 
Remote Sensing, IGARSS, FPL, FCCM). A estratégia de busca utiliza 
combinações de termos-chave incluindo ``hyperspectral'', ``FPGA'', 
``real-time'', ``embedded'', ``classification'', ``compression'', garantindo 
cobertura do estado da arte. Cada trabalho identificado é analisado segundo 
protocolo que extrai informações sobre arquitetura utilizada, técnicas 
implementadas, datasets de validação, métricas de performance, e limitações 
reportadas.

A catalogação de otimizações GPU segue metodologia similar, com ênfase em 
técnicas para plataformas embarcadas com restrições energéticas. O protocolo 
de análise inclui identificação de otimizações CUDA/OpenCL, estratégias de 
gerenciamento de memória, técnicas de paralelização, uso de precisão mista, e 
implementação de algoritmos de aprendizado de máquina. Trabalhos como 
\textcite{Diaz2019}, que implementa compressão hiperespectral em tempo real em 
GPUs embarcadas, e \textcite{Vaithianathan2025}, que discute desafios em 
arquiteturas de computação heterogênea, são analisados para extração de 
padrões de otimização e identificação de práticas.

A análise de trabalhos correlatos estende-se além de FPGA e GPU para incluir 
processadores ARM, DSPs, e arquiteturas heterogêneas. Esta análise permite 
identificação de tendências tecnológicas, lacunas na literatura, e 
oportunidades de integração entre plataformas. Trabalhos que abordam aspectos 
de integração de sistemas, como \textcite{Zhang2023} que apresenta revisão 
sobre classificação hiperespectral em UAV, recebem atenção para compreensão 
dos desafios de integração e validação em cenários reais.

A taxonomia de abordagens identificadas organiza as técnicas segundo critérios 
de classificação que incluem tipo de plataforma (FPGA, GPU, CPU, híbrida), 
estágio de processamento (pré-processamento, processamento principal, 
pós-processamento), técnica principal (compressão, classificação, redução 
dimensional), e métrica de otimização (throughput, latência, energia, 
precisão). Esta taxonomia facilita identificação de padrões, comparação entre 
abordagens, e síntese de recomendações.

\subsection{Modelagem Teórica Comparativa}
\label{subsec:modelagem_teorica}

Os modelos matemáticos de desempenho FPGA baseiam-se em análise de 
complexidade computacional e mapeamento para recursos de hardware disponíveis. 
Para cada técnica catalogada, desenvolve-se modelo que relaciona parâmetros de 
entrada (resolução espacial, número de bandas espectrais, taxa de frames) com 
utilização de recursos (DSPs, BRAMs, LUTs, flip-flops) e métricas de 
performance (throughput, latência, consumo energético). Os modelos consideram 
características específicas das famílias Zynq e UltraScale+, incluindo 
arquitetura de interconexão, hierarquia de memória, e capacidades de 
processamento em ponto fixo e flutuante.

A modelagem de throughput e latência GPU utiliza modelo de execução CUDA que 
considera características da arquitetura (número de SMs, cores CUDA, tensor 
cores), hierarquia de memória (registradores, shared memory, L1/L2 cache, 
memória global), e padrões de acesso a dados. Para cada algoritmo analisado, 
calcula-se ocupancy teórica, throughput de memória requerido, e latência de 
execução considerando overhead de comunicação CPU-GPU. Os modelos são 
calibrados com dados experimentais disponíveis na literatura para garantir 
precisão das estimativas.

A análise de complexidade computacional estabelece bounds teóricos para 
performance de diferentes algoritmos, independentemente da plataforma de 
implementação. Esta análise inclui complexidade temporal (Big-O notation), 
complexidade espacial (requisitos de memória), complexidade de comunicação 
(transferências de dados), e complexidade energética (operações por joule). 
A comparação de complexidades permite identificação de algoritmos 
intrinsecamente mais eficientes e estabelecimento de limites teóricos para 
performance.

A modelagem de consumo energético integra modelos específicos de cada 
plataforma com características dos algoritmos implementados. Para FPGA, 
utiliza-se modelo que considera consumo estático (leakage), consumo dinâmico 
(switching), e consumo de I/O. Para GPU, considera-se consumo base (idle), 
consumo computacional (proporcional à utilização), e consumo de memória 
(proporcional à largura de banda utilizada). Os modelos são validados com 
medições experimentais reportadas na literatura.

\subsection{Framework de Simulação Integrado}
\label{subsec:framework_simulacao}

O ambiente de simulação multi-plataforma integra simuladores específicos para 
cada tipo de arquitetura, permitindo avaliação comparativa sob condições 
controladas e padronizadas. Para FPGA, utiliza-se Vivado HLS para síntese de 
alto nível e Vivado para implementação, permitindo estimativa precisa de 
utilização de recursos e performance. Para GPU, emprega-se nsight Compute para 
profiling detalhado e simuladores de arquitetura para análise de performance 
teórica.

Os simuladores específicos incluem Vivado Design Suite para FPGA Xilinx, 
fornecendo estimativas de timing, power, e utilização de recursos após síntese 
e place-and-route. Para GPU NVIDIA, utiliza-se CUDA toolkit com profilers 
integrados (nsight Systems, nsight Compute) para análise detalhada de 
performance e identificação de gargalos. Simuladores de terceiros, como 
gem5-gpu para análise arquitetural detalhada, complementam as ferramentas 
proprietárias.

Os modelos de workload hiperespectral baseiam-se em caracterização detalhada 
de datasets públicos (AVIRIS, Pavia, Salinas) e aplicações típicas 
(classificação, detecção de anomalias, unmixing espectral). Cada workload é 
caracterizado por distribuição de intensidades espectrais, padrões espaciais, 
requisitos computacionais, e sensibilidade a diferentes tipos de otimização. 
Esta caracterização permite simulação realística de diferentes cenários de 
aplicação.

As métricas de comparação padronizadas incluem throughput (fps, MB/s), latência 
(ms end-to-end, ms por estágio), utilização de recursos (percentual de DSPs, 
BRAMs, cores GPU), consumo energético (W médio, J por frame), e qualidade de 
resultado (accuracy, F1-score, PSNR). Todas as métricas são medidas sob 
condições padronizadas para garantir comparabilidade entre diferentes técnicas 
e plataformas.

\section{Estudo Comparativo de Arquiteturas}
\label{sec:estudo_comparativo}

A análise sistemática e comparativa das diferentes abordagens arquiteturais 
estrutura-se em três dimensões principais: técnicas específicas para FPGA 
Xilinx, otimizações para GPU embarcadas, e estudo comparativo cross-platform 
que identifica trade-offs fundamentais entre as arquiteturas. Esta seção 
apresenta metodologia detalhada para cada dimensão, estabelecendo critérios 
objetivos para comparação e frameworks analíticos para identificação de 
soluções ótimas.

\subsection{Análise de Técnicas FPGA Xilinx}\label{subsec:analise_fpga}

A comparação entre High-Level Synthesis (HLS) e HDL tradicional baseia-se em 
análise quantitativa de produtividade de desenvolvimento, qualidade de 
resultados, e flexibilidade de otimização. Para cada técnica de processamento 
hiperespectral identificada na literatura, avalia-se a adequação de HLS versus 
HDL considerando complexidade do algoritmo, requisitos de performance, e 
restrições de recursos. O trabalho de \textcite{Hwang2011} propõe metodologia de design consciente de 
energia para compressão de imagens hiperespectrais, demonstrando a importância 
de considerações energéticas no desenvolvimento de sistemas embarcados. 
\textcite{Khaki2025} complementam esta análise demonstrando que FPGAs 
alcançam speedups de 3.31× versus CPU para MobileNetV2, com quantização INT8 
preservando acurácia enquanto proporciona ganhos de eficiência. Algoritmos 
com padrões de acesso irregulares podem requerer otimizações específicas 
dependendo da arquitetura de implementação.

A análise de paralelização em pipeline versus paralelo examina diferentes 
estratégias de mapeamento de algoritmos para arquitetura FPGA. A paralelização 
em pipeline é avaliada para algoritmos com dependências sequenciais, como 
filtros IIR e transformadas, considerando throughput máximo, latência inicial, 
e utilização de recursos. A paralelização espacial é analisada para operações 
independentes, como convolução e classificação pixel-wise, avaliando 
escalabilidade, eficiência de recursos, e balanceamento de carga. A comparação 
utiliza métricas de eficiência computacional (operações por LUT por segundo) e 
eficiência energética (operações por watt).

As otimizações de memória e DSP constituem aspecto crítico para performance em 
aplicações hiperespectrais devido ao grande volume de dados e intensidade 
computacional. A análise inclui estratégias de particionamento de memória 
(on-chip BRAM, UltraRAM, memória externa), técnicas de buffering para 
maximização de reuso de dados, e otimizações específicas de DSP para operações 
de ponto fixo e flutuante. Estudos na literatura demonstram que otimizações de memória podem resultar em 
melhorias significativas em throughput através de redução de acessos à memória 
externa, conforme evidenciado por trabalhos como \textcite{ArucuIliev2025} que 
apresenta análise comparativa de eficiência energética em processamento 
hiperespectral acelerado por FPGA.

As técnicas de particionamento de dados examinam diferentes estratégias para 
divisão de imagens hiperespectrais considerando características espaciais e 
espectrais. O particionamento espacial (tiles, strips, blocks) é avaliado 
quanto a overhead de comunicação, balanceamento de carga, e requisitos de 
memória. O particionamento espectral (bandas individuais, grupos de bandas, 
PCA) é analisado considerando preservação de informação, redução de 
complexidade, e adequação a diferentes algoritmos de processamento.

\subsection{Análise de Técnicas GPU Embarcadas}
\label{subsec:analise_gpu}

A comparação entre otimizações CUDA e OpenCL para GPU embarcadas considera 
performance, portabilidade, e facilidade de desenvolvimento. CUDA oferece 
otimizações específicas para arquiteturas NVIDIA e ferramentas de profiling 
avançadas, enquanto OpenCL proporciona portabilidade entre diferentes 
fabricantes com potencial perda de performance. A análise quantitativa baseia-se 
em implementações de referência de algoritmos hiperespectrais típicos 
(classificação SVM, PCA, unmixing espectral) em ambas as plataformas, 
comparando throughput, utilização de recursos, e consumo energético.

As estratégias de memory coalescing e shared memory são analisadas 
considerando padrões de acesso típicos em processamento hiperespectral. 
Imagens hiperespectrais apresentam localidade espacial (pixels vizinhos) e 
espectral (bandas adjacentes), permitindo otimizações específicas de acesso à 
memória. A análise inclui técnicas de tiling para maximização de coalescing, 
uso de shared memory para reuso de dados entre threads de um warp, e 
estratégias de prefetching para ocultação de latência de memória.

A ocupancy optimization examina técnicas para maximização da utilização de 
recursos computacionais considerando limitações específicas de GPU embarcadas. 
Estas plataformas possuem menor número de SMs e menor largura de banda de 
memória comparadas a GPU de desktop, requerendo estratégias específicas de 
otimização. A análise inclui balanceamento entre paralelismo e utilização de 
memória, técnicas de loop unrolling e tiling, e estratégias de scheduling para 
maximização de throughput.

As técnicas de mixed precision e quantização são fundamentais para aplicações 
embarcadas devido a restrições de memória e largura de banda. A análise 
examina impacto da redução de precisão (FP32 para FP16, INT8) na qualidade de 
resultados para diferentes algoritmos hiperespectrais. Pesquisas na área demonstram que técnicas de quantização podem proporcionar 
reduções significativas no consumo de memória e melhorias no throughput, 
mantendo níveis aceitáveis de precisão. O trabalho de \textcite{Diaz2019} 
explora otimizações para compressão hiperespectral em tempo real, demonstrando 
a viabilidade de implementações eficientes em GPUs embarcadas. 
\textcite{Khaki2025} demonstram que sistemas heterogêneos utilizando 
frameworks como Vitis AI permitem deployment escalável de CNNs, com FPGAs 
sendo adequadas para edge computing com limitações energéticas, enquanto GPUs 
mantêm vantagem em velocidade absoluta com latências de 13.64s versus 100.4s 
em FPGA para MobileNetV2.

\subsection{Estudo Comparativo Cross-Platform}
\label{subsec:comparativo_cross_platform}

A análise de trade-offs fundamentais entre FPGA e GPU estabelece framework 
quantitativo para comparação considerando diferentes métricas de performance e 
cenários de aplicação. FPGA oferece baixa latência, determinismo temporal, e 
alta eficiência energética, sendo adequada para aplicações com requisitos 
rigorosos de tempo real. GPU proporciona alto throughput, flexibilidade de 
programação, e capacidade de processamento de grandes volumes de dados, sendo 
adequada para aplicações batch ou com tolerância a latência variável.

A análise de adequação por workload categoriza diferentes tipos de 
processamento hiperespectral segundo características computacionais e 
adequação a cada arquitetura. Algoritmos com paralelismo regular e alta 
intensidade aritmética (convolução, transformadas) favorecem GPU, enquanto 
algoritmos com requisitos de baixa latência e padrões de acesso irregulares 
(classificação online, controle adaptativo) favorecem FPGA. A análise 
quantitativa baseia-se em implementações de referência e modelos teóricos de 
performance.

As métricas de eficiência energética constituem critério fundamental para 
aplicações embarcadas, especialmente em UAV com restrições de autonomia. A 
análise compara consumo energético por operação (joules por classificação, 
joules por frame processado) entre diferentes arquiteturas, considerando 
consumo estático e dinâmico. O trabalho de \textcite{Hwang2011} propõe metodologia de design consciente de 
energia para compressão de imagens hiperespectrais, destacando a importância 
da eficiência energética em sistemas embarcados. O estudo de \textcite{ArucuIliev2025} 
apresenta análise comparativa sobre eficiência energética em processamento 
hiperespectral acelerado por FPGA, enquanto \textcite{Diaz2019} demonstra 
implementação eficiente em GPUs embarcadas para compressão em tempo real.

A análise de escalabilidade e flexibilidade examina capacidade de adaptação 
das arquiteturas a diferentes requisitos de aplicação. FPGA oferece 
flexibilidade arquitetural através de reconfiguração, permitindo otimização 
específica para cada algoritmo, mas requer tempo de desenvolvimento 
significativo. GPU proporciona flexibilidade de software através de 
programação de alto nível, permitindo adaptação rápida a novos algoritmos, mas 
com arquitetura fixa que pode não ser ótima para todos os casos.

\section{Protocolo de Simulação e Benchmarking}
\label{sec:protocolo_simulacao}

A metodologia detalhada para simulação e avaliação de desempenho estabelece 
ambiente controlado para comparação objetiva entre diferentes técnicas e 
arquiteturas. O protocolo estrutura-se em três componentes principais: 
configuração de ambiente de simulação padronizado, definição de datasets e 
cargas de trabalho representativas, e estabelecimento de protocolo de medição 
e análise estatística. Esta padronização garante reprodutibilidade dos 
resultados e comparabilidade entre diferentes estudos.

\subsection{Ambiente de Simulação}\label{subsec:ambiente_simulacao}

A configuração de ferramentas para FPGA baseia-se em Vivado Design Suite 
2023.1 com Vitis HLS para síntese de alto nível, proporcionando ambiente 
completo para desenvolvimento, simulação e análise de performance. As 
configurações incluem targets específicos para dispositivos Zynq UltraScale+ 
(XCZU7EV, XCZU9EG) com caracterização detalhada de recursos disponíveis, 
frequências de operação, e capacidades de DSP. O ambiente é configurado para 
simulação funcional, síntese lógica, e implementação física, permitindo 
estimativas precisas de timing, power, e utilização de recursos.

O setup de simuladores GPU utiliza CUDA Toolkit 12.0 com nsight Compute e 
nsight Systems para profiling detalhado de aplicações CUDA. As configurações 
incluem targets para dispositivos Jetson (TX2, Xavier, Orin) com modelagem 
precisa de arquitetura (número de SMs, cores CUDA, tensor cores), hierarquia 
de memória (capacidade e largura de banda), e características térmicas. 
Simuladores complementares, como gem5-gpu, proporcionam análise arquitetural 
detalhada para exploração de design space.

Os modelos de hardware sintéticos complementam simuladores comerciais para 
exploração de configurações não disponíveis comercialmente ou análise de 
tendências futuras. Estes modelos baseiam-se em extrapolação de características 
de dispositivos existentes, considerando roadmaps tecnológicos e limitações 
físicas fundamentais. A validação dos modelos sintéticos utiliza comparação 
com dispositivos reais para garantir precisão das estimativas.

A validação dos ambientes de simulação estabelece protocolos de verificação 
que comparam resultados simulados com medições experimentais disponíveis na 
literatura. Para cada ferramenta de simulação, define-se conjunto de casos de 
teste com resultados conhecidos, estabelecendo intervalos de confiança para 
diferentes tipos de estimativa. Esta validação é essencial para garantir 
credibilidade dos resultados teóricos propostos.

\subsection{Datasets e Cargas de Trabalho}\label{subsec:datasets_workloads}

A caracterização de workloads hiperespectrais baseia-se em análise estatística 
detalhada de datasets públicos representativos de diferentes aplicações. O 
dataset AVIRIS Indian Pines (145×145×200 bandas) representa aplicações 
agrícolas com 16 classes de cobertura do solo, caracterizado por alta 
variabilidade espectral e desbalanceamento entre classes. O dataset Pavia 
University (610×340×103 bandas) representa aplicações urbanas com 9 classes, 
caracterizado por alta resolução espacial e contraste espectral. O dataset 
Salinas Valley (512×217×204 bandas) representa agricultura de precisão com 16 
classes de culturas, caracterizado por sutis diferenças espectrais entre 
classes similares.

Os datasets sintéticos para benchmarking são gerados através de modelos 
estatísticos que preservam características espectrais e espaciais dos dados 
reais, permitindo controle preciso de parâmetros como resolução, número de 
bandas, nível de ruído, e complexidade de classificação. Esta abordagem 
permite avaliação sistemática de sensibilidade dos algoritmos a diferentes 
condições operacionais e identificação de limitações específicas.

As métricas de complexidade computacional caracterizam cada workload segundo 
intensidade aritmética (operações por pixel), intensidade de memória (bytes 
por operação), padrões de acesso a dados (sequencial, aleatório, localidade), 
e paralelismo disponível (independência entre operações). Esta caracterização 
permite predição de adequação a diferentes arquiteturas e identificação de 
gargalos potenciais.

Os casos de teste representativos incluem cenários típicos de aplicação com 
requisitos específicos de performance e qualidade. O cenário de agricultura de 
precisão simula UAV a 100m de altitude processando área de 1 hectare em tempo 
real para detecção de estresse hídrico. O cenário de monitoramento urbano 
simula classificação LULC contínua a 30 fps para análise de mudanças. O 
cenário de aplicação de emergência simula detecção de incêndios florestais com 
latência crítica inferior a 100ms.

\subsection{Protocolo de Medição e Análise}
\label{subsec:protocolo_medicao}

As métricas de desempenho incluem throughput (fps, MB/s, pixels/s), latência 
(ms end-to-end, ms por estágio, percentil 95\%), e variabilidade temporal 
(desvio padrão, jitter). O throughput é medido sob condições de carga 
sustentada para identificar performance máxima e degradação temporal. A 
latência é medida desde aquisição da imagem até disponibilização do resultado, 
incluindo overhead de comunicação e sincronização. A variabilidade temporal é 
crítica para aplicações de tempo real que requerem comportamento determinístico.

A análise de utilização de recursos quantifica eficiência de mapeamento de 
algoritmos para hardware disponível. Para FPGA, mede-se utilização de LUTs, 
flip-flops, DSPs, e BRAMs, identificando recursos limitantes e oportunidades 
de otimização. Para GPU, mede-se ocupancy de SMs, utilização de memória, e 
largura de banda consumida, identificando gargalos e possibilidades de 
melhoria.

O profiling de consumo energético utiliza modelos validados para estimativa 
de consumo estático e dinâmico de cada arquitetura. As medições incluem 
potência média, potência pico, energia por frame processado, e distribuição de 
consumo entre diferentes componentes (computação, memória, I/O). Esta análise 
é fundamental para aplicações embarcadas com restrições energéticas rigorosas.

A análise estatística comparativa emprega testes de significância para 
validação de diferenças de performance entre técnicas e arquiteturas. A 
metodologia inclui análise de variância (ANOVA) para comparação múltipla, 
testes post-hoc para identificação de diferenças específicas, e intervalos de 
confiança para quantificação de incerteza. Esta rigor estatístico garante 
validade científica das conclusões propostas.

\section{Metodologia de Síntese Arquitetural}
\label{sec:sintese_arquitetural}

O processo de proposição da arquitetura otimizada baseia-se nos resultados 
comparativos obtidos nas etapas anteriores, utilizando metodologia sistemática 
de tomada de decisão multi-critério. Esta seção apresenta a análise 
multi-critério para ponderação de diferentes métricas de performance, o 
processo de design space exploration para identificação de configurações 
ótimas, e a metodologia de síntese para composição da arquitetura final 
proposta.

\subsection{Análise Multi-Critério}\label{subsec:analise_multicriterio}

A ponderação de métricas de desempenho utiliza metodologia AHP (Analytic 
Hierarchy Process) para estabelecimento de pesos relativos entre diferentes 
critérios segundo cenários específicos de aplicação. Para aplicações de 
agricultura de precisão, prioriza-se eficiência energética e autonomia 
operacional, atribuindo pesos elevados a consumo energético e throughput 
sustentado. Para aplicações de monitoramento de emergência, prioriza-se baixa 
latência e confiabilidade, enfatizando tempo de resposta e robustez a falhas.

A análise de Pareto para trade-offs identifica conjunto de soluções 
não-dominadas no espaço multidimensional de objetivos conflitantes. As 
dimensões incluem performance (throughput, latência), eficiência (energia, 
recursos), qualidade (precisão, robustez), e custo (desenvolvimento, hardware). 
A fronteira de Pareto estabelece limites teóricos para otimização simultânea 
e identifica regiões do espaço de soluções com melhor compromisso entre 
objetivos conflitantes.

O método AHP para tomada de decisão estrutura-se em hierarquia de critérios 
com comparações pareadas para estabelecimento de pesos relativos. O primeiro 
nível inclui critérios principais (performance, eficiência, qualidade, custo), 
o segundo nível decompõe cada critério em sub-critérios específicos, e o 
terceiro nível avalia alternativas técnicas segundo cada sub-critério. A 
metodologia inclui verificação de consistência das comparações e análise de 
sensibilidade para validação da robustez das decisões.

O ranking de técnicas por cenário utiliza pontuação ponderada baseada nos 
pesos AHP e performance normalizada de cada técnica segundo diferentes 
critérios. Esta abordagem permite identificação de técnicas mais adequadas a 
cenários específicos e composição de arquiteturas híbridas que combinam 
diferentes técnicas segundo suas vantagens comparativas.

\subsection{Design Space Exploration}\label{subsec:design_space}

O mapeamento do espaço de soluções estrutura-se em dimensões que incluem tipo 
de arquitetura (FPGA, GPU, híbrida), técnicas de processamento (compressão, 
classificação, fusão), configurações de hardware (frequência, memória, 
paralelismo), e parâmetros algorítmicos (precisão, janela, threshold). Cada 
ponto no espaço representa configuração específica com performance estimada 
segundo modelos teóricos desenvolvidos.

A identificação de configurações ótimas utiliza algoritmos de otimização 
multi-objetivo (NSGA-II, MOPSO) para exploração eficiente do espaço de 
soluções. Os algoritmos de otimização multi-objetivo consideram restrições operacionais (consumo máximo, 
latência máxima, precisão mínima) e identificam conjunto de configurações 
Pareto-ótimas que representam melhores compromissos entre objetivos 
conflitantes.

A análise de sensibilidade paramétrica examina robustez das configurações 
ótimas a variações em parâmetros de entrada e condições operacionais. Esta 
análise inclui sensibilidade a variações em datasets (resolução, ruído, 
complexidade), condições ambientais (temperatura, vibração), e degradação de 
hardware (aging, variabilidade de processo). Configurações com baixa 
sensibilidade são preferidas por proporcionarem performance mais previsível.

A otimização multi-objetivo considera simultaneamente múltiplas funções 
objetivo conflitantes, utilizando técnicas de programação matemática e 
algoritmos evolutivos. A formulação matemática inclui funções objetivo para 
maximização de throughput, minimização de latência, minimização de consumo 
energético, e maximização de precisão, sujeitas a restrições de recursos 
disponíveis e requisitos operacionais.

\subsection{Síntese da Arquitetura Proposta}
\label{subsec:sintese_arquitetura}

A metodologia de composição arquitetural baseia-se em princípios de engenharia 
de sistemas que consideram interfaces, dependências, e otimização global. A 
arquitetura proposta integra componentes FPGA e GPU segundo pipeline otimizado 
que explora vantagens comparativas de cada plataforma: FPGA para 
pré-processamento de baixa latência e alta eficiência, GPU para processamento 
principal de alto throughput.

A integração de técnicas complementares utiliza análise de compatibilidade e 
sinergia entre diferentes abordagens identificadas na literatura. Técnicas de 
compressive sensing em FPGA são integradas com decompressão CUDA em GPU, 
técnicas de seleção de bandas são combinadas com classificação otimizada, e 
estratégias de fusão multi-sensor complementam processamento individual de 
cada modalidade.

A especificação da solução otimizada inclui arquitetura detalhada com 
componentes específicos, interfaces de comunicação, protocolos de 
sincronização, e configurações de software. A especificação é suficientemente 
detalhada para permitir implementação futura, incluindo diagramas de blocos, 
especificações de timing, estimativas de recursos, e análise de viabilidade 
técnica.

A validação teórica da proposta utiliza modelos matemáticos desenvolvidos para 
estimativa de performance da arquitetura integrada. A validação inclui análise 
de throughput end-to-end considerando gargalos e overhead de comunicação, 
estimativa de consumo energético total, análise de latência incluindo 
sincronização entre componentes, e avaliação de robustez a variações 
operacionais.

\section{Planejamento de Execução}\label{sec:planejamento_execucao}

O cronograma e organização temporal da pesquisa comparativa estrutura-se em 
quatro fases sequenciais com marcos específicos e entregáveis bem definidos. 
Esta seção apresenta o planejamento detalhado de cada fase, incluindo 
atividades, recursos necessários, marcos intermediários, e estratégias de 
gestão de riscos para garantir execução dentro do prazo estabelecido.

\subsection{Fases de Desenvolvimento}\label{subsec:fases_desenvolvimento}

A Fase 1 de Levantamento e Modelagem, com duração de 3 meses, concentra-se na 
consolidação da base teórica e desenvolvimento dos modelos matemáticos. As 
atividades incluem revisão sistemática complementar para identificação de 
trabalhos recentes, catalogação detalhada de técnicas segundo taxonomia 
proposta, desenvolvimento de modelos teóricos de performance para cada 
plataforma, e validação inicial dos modelos com dados experimentais da 
literatura. Esta fase estabelece fundação sólida para as análises comparativas 
subsequentes.

A Fase 2 de Simulação Comparativa, com duração de 4 meses, implementa o 
framework de simulação e executa análises comparativas sistemáticas. As 
atividades incluem configuração de ambientes de simulação (Vivado, CUDA), 
implementação de modelos de workload hiperespectral, execução de simulações 
para diferentes configurações e cenários, coleta e análise estatística de 
resultados, e identificação de padrões e tendências. Esta fase representa o 
núcleo da contribuição metodológica da pesquisa.

A Fase 3 de Análise e Síntese, com duração de 3 meses, processa os resultados 
das simulações para proposição da arquitetura otimizada. As atividades incluem 
análise multi-critério dos resultados comparativos, design space exploration 
para identificação de configurações ótimas, síntese da arquitetura proposta 
baseada nos resultados, validação teórica da proposta, e desenvolvimento de 
recomendações para diferentes cenários de aplicação.

A Fase 4 de Consolidação e Redação, com duração de 2 meses, finaliza a 
documentação da pesquisa e prepara materiais para defesa. As atividades 
incluem redação dos capítulos de resultados e conclusões, revisão e 
consolidação de todo o documento, preparação de material de apresentação, 
revisão final por orientador e colaboradores, e submissão da versão final da 
dissertação.

\subsection{Marcos e Entregáveis}\label{subsec:marcos_entregaveis}

O Marco M1 - Catálogo de Técnicas e Modelos marca a conclusão da Fase 1 com 
entregáveis que incluem catálogo sistemático de técnicas FPGA e GPU, taxonomia 
de abordagens identificadas, modelos matemáticos validados para cada 
plataforma, e relatório de fundamentação teórica. Este marco estabelece base 
sólida para as fases subsequentes e permite validação da metodologia proposta.

O Marco M2 - Resultados de Simulação Comparativa marca a conclusão da Fase 2 
com entregáveis que incluem framework de simulação implementado e validado, 
base de dados completa com resultados de simulação, análise estatística 
comparativa entre técnicas e arquiteturas, e relatório de resultados 
experimentais. Este marco representa a principal contribuição empírica da 
pesquisa.

O Marco M3 - Arquitetura Proposta e Validada marca a conclusão da Fase 3 com 
entregáveis que incluem especificação detalhada da arquitetura otimizada, 
análise de trade-offs e justificativa das escolhas, validação teórica da 
proposta com estimativas de performance, e recomendações para implementação 
futura. Este marco consolida a contribuição teórica principal da dissertação.

O Marco M4 - Dissertação Completa marca a conclusão da Fase 4 com entregáveis 
que incluem documento final da dissertação revisado e aprovado, material de 
apresentação para defesa, artigos científicos preparados para submissão, e 
código fonte do framework de simulação documentado. Este marco finaliza todos 
os produtos da pesquisa.

\subsection{Gestão de Riscos e Contingências}
\label{subsec:gestao_riscos}

Os riscos de limitação de simuladores incluem possível indisponibilidade de 
licenças de software proprietário, limitações de precisão dos modelos de 
simulação, e incompatibilidade entre diferentes versões de ferramentas. As 
estratégias de mitigação incluem identificação de alternativas open-source 
(GHDL, Verilator para FPGA; OpenCL para GPU), desenvolvimento de modelos 
próprios quando necessário, e manutenção de múltiplas versões de ferramentas 
para garantir compatibilidade.

As contingências para validação de modelos incluem possível divergência entre 
resultados simulados e dados experimentais da literatura, limitações de 
datasets públicos disponíveis, e dificuldades de reprodução de resultados de 
trabalhos anteriores. As estratégias incluem desenvolvimento de múltiplos 
modelos com diferentes níveis de abstração, utilização de datasets sintéticos 
com características controladas, e colaboração com autores de trabalhos 
relevantes para esclarecimento de detalhes de implementação.

As alternativas metodológicas incluem possível necessidade de ajustes na 
abordagem comparativa, limitações de escopo devido a restrições de tempo ou 
recursos, e mudanças nos objetivos baseadas em resultados intermediários. As 
estratégias incluem definição de objetivos mínimos e estendidos para cada fase, 
flexibilidade na seleção de técnicas e plataformas analisadas, e adaptação da 
metodologia baseada em lições aprendidas.

O plano de mitigação de riscos estabelece protocolos de monitoramento contínuo 
do progresso, identificação precoce de problemas potenciais, e ativação de 
estratégias de contingência quando necessário. O plano inclui reuniões 
semanais de acompanhamento, relatórios mensais de progresso, e revisões 
trimestrais de escopo e cronograma para garantir execução bem-sucedida da 
pesquisa.

\section{Considerações Finais do Capítulo}
\label{sec:consideracoes_finais_metodologia}

A metodologia proposta estabelece framework sistemático para investigação 
comparativa de técnicas de processamento hiperespectral em arquiteturas 
heterogêneas, preenchendo lacuna metodológica identificada na literatura. A 
abordagem exploratória-comparativa baseada em simulação e modelagem teórica 
permite análise abrangente do espaço de soluções disponíveis, identificação de 
trade-offs fundamentais entre diferentes arquiteturas, e proposição de 
arquitetura otimizada baseada em critérios objetivos e quantitativos.

A justificativa das escolhas metodológicas baseia-se na análise crítica do 
estado da arte que revela ausência de estudos sistemáticos comparativos e 
predominância de trabalhos focados em técnicas isoladas. A metodologia proposta 
aborda estas limitações através de framework unificado que permite comparação 
objetiva entre diferentes abordagens, considerando múltiplos critérios de 
performance e diferentes cenários de aplicação. A ênfase em simulação e 
modelagem teórica justifica-se pela complexidade e custo de implementação 
experimental em múltiplas plataformas heterogêneas.

As perspectivas de contribuição do estudo comparativo para o estado da arte 
incluem estabelecimento de metodologia padronizada para comparação de técnicas 
hiperespectrais, identificação quantitativa de trade-offs entre diferentes 
arquiteturas, proposição de arquitetura otimizada baseada em análise 
sistemática, e desenvolvimento de recomendações práticas para seleção de 
técnicas segundo requisitos específicos de aplicação. Estas contribuições 
representam avanço significativo no conhecimento científico da área e 
estabelecem fundação para desenvolvimentos futuros.

A metodologia apresentada neste capítulo estabelece base sólida para execução 
da pesquisa proposta, garantindo rigor científico, reprodutibilidade dos 
resultados, e relevância prática das contribuições. A estruturação em fases 
bem definidas com marcos específicos permite acompanhamento do progresso e 
ajustes quando necessário, enquanto as estratégias de gestão de riscos 
asseguram execução bem-sucedida dentro do prazo estabelecido.
